\section{Apéndice: plantillas de Prompt y de Tool}

\subsection{Arquitectura del Prompt Template}

El prompt template que recomienda Burak consta de tres partes:

\begin{enumerate}
    \item Context: incluye los antecedentes de la tarea, las restricciones de policy y la intención del usuario
    \item Instruction: especifica los requisitos de salida del paso actual (por ejemplo, «enumera las top-3 soluciones»)
    \item Tool Usage: enumera los métodos de invocación de las herramientas disponibles, junto con un success indicator
\end{enumerate}

\begin{knowledgebox}{Lista de verificación del Prompt Template}
Cada vez que se actualiza el prompt, hay que verificar:
1. Si el Context se actualizó (policy / data)
2. Si la Instruction sigue siendo clara
3. Si las Tools \& examples son reproducibles
\end{knowledgebox}

\subsection{Ejemplo de Tool Manifest}

Un tool manifest típico incluye los siguientes campos:

\begin{itemize}
    \item tool name
    \item description
    \item input schema (tipos de parámetro, valores posibles)
    \item latency range
    \item safety constraints (qué user group no puede usarla)
\end{itemize}

Incluso si la tool es solo una API interna, se recomienda redactarla como manifest; el equipo de logística puede usarla como manual de operación.

\subsection{Sistema de etiquetas Failure Tag}

Burak recomienda unificar la failure tag taxonomy:

\begin{itemize}
    \item hallucination
    \item tool failure
    \item context expiration
    \item safety violation
    \item latency breach
\end{itemize}

Cada tag está asociado a un flujo de procesamiento: por ejemplo, `tool failure` activa fallback or escalate, y `hallucination` activa LLM retrace path.

\subsection{Resumen de la sección}

Las plantillas de prompt y de tool son la base de agent reproducibility; una vez estructurados los modos de fallo, el producto y la ingeniería pueden corregirlos con iteraciones más pequeñas.

